\section{Cost and Performance Model}
\label{sec:cost}

Because the fabric is driven entirely over the serial link, its externally
observed performance is set not by the number of PEs but by the number and size
of the protocol transactions a kernel issues. This section gives a first-order,
purely analytical model of that cost, derived from the protocol of
Section~\ref{sec:proto}; it is a wire-level accounting, not a measurement, and
deliberately ignores the negligible on-device compute time (a \cmd{RUN} of $N$
steps is $N$ cycles at $100\,\mathrm{MHz}$, dwarfed by the serial exchange).

\subsection{Per-command wire cost}

Every transaction is a round trip: the host sends a command and its payload and
waits for the reply. Table~\ref{tab:cost} lists the byte counts that follow
directly from the payload and reply sizes of Table~\ref{tab:proto}, for the
$4\times4$, 16-bit fabric.

\begin{table}[t]
  \centering
  \caption{Wire cost per transaction (bytes each way).}
  \label{tab:cost}
  \small
  \begin{tabular}{@{}lrrl@{}}
    \toprule
    Cmd & TX & RX & Note \\
    \midrule
    \cmd{ID}  & 1  & 5  & signature $+$ geometry \\
    \cmd{CFG} & 66 & 1  & $16\times4$ words $+$ cmd $+$ cksum \\
    \cmd{WR}  & 18 & 1  & $8\times2$ words $+$ cmd $+$ cksum \\
    \cmd{RUN} & 2  & 1  & cmd $+$ step count \\
    \cmd{RD}  & 1  & 33 & $16\times2$ regs $+$ cksum \\
    \cmd{RST} & 1  & 1  & clears registers \\
    \bottomrule
  \end{tabular}
\end{table}

\subsection{Per-kernel transaction count}

Composing these primitives as the software stack does (Section~\ref{sec:sw})
gives the cost of each kernel. Writing $\lceil\cdot\rceil$ for the ceiling and
$T_C=\lceil N/\Cols\rceil$, $T_R=\lceil M/\Rows\rceil$ for the number of column
and row tiles, the transaction counts are:

\begin{itemize}
  \item \emph{Element-wise}, $n$ elements ($\Rows$ lanes per pass):
    one \cmd{CFG}, then $\lceil n/\Rows\rceil$ rounds of
    \cmd{WR}$+$\cmd{RUN}$+$\cmd{RD}, i.e.\ $1+3\lceil n/\Rows\rceil$
    transactions.
  \item \emph{Reduction} (dot product), $n$ elements: one \cmd{CFG}, one
    \cmd{RST}, then $n$ rounds of \cmd{WR}$+$\cmd{RUN}, and one final
    \cmd{RD}: $3+2n$ transactions.
  \item \emph{Systolic} $M\times N$ matrix--vector: one \cmd{CFG} per column
    tile and, per tile pair, one \cmd{RST}, $\Rows$ \cmd{WR}$+$\cmd{RUN}
    rounds and one \cmd{RD}, i.e.\ $T_C\,(1+10\,T_R)$ transactions.
\end{itemize}

The element-wise cost is dominated by the $\cmd{RD}$ reply (33~bytes) and is
amortised only four ways, so throughput is essentially flat in the array size:
a wider fabric would inject and read more lanes per round but would not change
the per-round round-trip structure.

\subsection{Latency model}

At $115200$~baud, 8N1, each byte occupies ten symbol times, so a byte costs
$t_b = 10/115200 \approx 86.8\,\mu\mathrm{s}$. Modelling a transaction as
the serial time of its two halves, an element-wise round
(\cmd{WR}$+$\cmd{RUN}$+$\cmd{RD} $=21$~TX $+35$~RX bytes) takes about
$56\,t_b \approx 4.9\,\mathrm{ms}$ for four results, i.e.\ roughly
$1.2\,\mathrm{ms}$ per element, independent of what the ALU computes. The
one-time \cmd{CFG} adds $67\,t_b \approx 5.8\,\mathrm{ms}$.

Two conclusions follow, and both point away from the array and toward the link.
First, the fixed per-round overhead---a full \cmd{RD} of all sixteen registers
regardless of how many lanes were used, plus command and checksum bytes---means
small transfers are latency-bound. Second, the natural optimisation is to
\emph{batch}: coalescing the write, run and read-back of many chunks into fewer,
larger transactions would raise the useful-payload fraction and cut the number
of round trips, addressing the true bottleneck identified in
Section~\ref{sec:concl} without touching the fabric.
